% ─────────────────────────────────────────────────
%  RÉSUMÉ
% ─────────────────────────────────────────────────
\chapter*{Résumé}
\phantomsection
\addcontentsline{toc}{chapter}{Résumé}
\markboth{Résumé}{Résumé}

\hspace{0.5cm}
Ce mémoire présente la conception et la réalisation des modules
\textbf{Relation Client} et \textbf{SGI} (Société de Gestion et
d'Intermédiation) de la plateforme \textbf{I-SIB}. Le travail s'inscrit dans le
besoin de centraliser la connaissance client et d'accompagner les opérations
d'investissement dans un système d'information bancaire modulaire.

\hspace{0.5cm}
La solution conçue repose sur une architecture microservices et une logique
\textit{full API}, afin de séparer les responsabilités métier et de faciliter
l'évolution progressive de la plateforme. Chaque service métier expose ses
fonctions via une passerelle API et conserve sa propre base de données, des
composants communs assurant le cache, la messagerie et la gestion des accès. La
plateforme a été déployée en conteneurs sur un VPS AlmaLinux. Pour la partie
détaillée, nous avons retenu la création d'un prospect comme scénario
représentatif du module Relation Client ; le module SGI, lui, est présenté à
travers les résultats fonctionnels liés à la consultation des informations
d'investissement et à la préparation des ordres.

\hspace{0.5cm}
Les résultats obtenus portent sur le suivi client, la prospection, la
consultation des informations d'investissement et la préparation des ordres. Les
cotations BRVM utilisées pour la validation proviennent d'une source publique
tierce qui relaie quotidiennement les données de la BRVM, faute d'un accès
direct au flux de marché officiel. Les
limites identifiées concernent notamment l'accès à un flux de marché officiel,
la consolidation des modèles d'intelligence artificielle, l'audit de sécurité et
la préparation de la production.

\vspace{0.8cm}
\noindent
\textbf{Mots-clés :} système d'information bancaire, relation client, SGI,
CRM, microservices, intelligence artificielle, BRVM, I-SIB.

\cleardoublepage

% ─────────────────────────────────────────────────
%  ABSTRACT
% ─────────────────────────────────────────────────
\chapter*{Abstract}
\phantomsection
\addcontentsline{toc}{chapter}{Abstract}
\markboth{Abstract}{Abstract}

\hspace{0.5cm}
This dissertation presents the design and implementation of the \textbf{Customer
Relationship} and \textbf{SGI} (Securities Management and Intermediation)
modules of the \textbf{I-SIB} platform. The work addresses the need to
centralize customer knowledge and support investment operations within a modular
banking information system.

\hspace{0.5cm}
The designed solution is based on a microservices architecture and a full API
approach, in order to separate business responsibilities and support the
progressive evolution of the platform. Each business service exposes its
functions through an API gateway and keeps its own database, with shared
components handling caching, messaging and access management. The platform was
deployed in containers on an AlmaLinux VPS. For the detailed part, we selected prospect creation
as a representative scenario of the Customer Relationship module; the SGI
module, in turn, is presented through the functional results related to
investment information consultation and order preparation.

\hspace{0.5cm}
The results cover customer monitoring, prospect tracking, investment information
consultation and order preparation. The BRVM quotes used for validation come
from a third-party public source that relays the BRVM data daily, without direct
access to the official market feed. The main limitations concern access to an official market
data feed, consolidation of artificial intelligence models, security auditing
and production readiness.

\vspace{0.8cm}
\noindent
\textbf{Keywords:} banking information system, customer relationship, SGI, CRM,
microservices, artificial intelligence, BRVM, I-SIB.

\cleardoublepage
